\documentclass[11pt,a4paper]{article}

\usepackage[utf8]{inputenc}
\usepackage[T1]{fontenc}
\usepackage[spanish]{babel}
\usepackage{amsmath,amssymb,mathtools}
\usepackage{geometry}
\usepackage{enumitem}
\usepackage{array}
\usepackage{lmodern}
\usepackage{microtype}
\usepackage{xcolor}

\geometry{margin=2.4cm}
\setlength{\parindent}{0pt}
\setlength{\parskip}{0.5em}

\newcommand{\R}{\mathbb{R}}
\newcommand{\K}{\mathbb{K}}
\newcommand{\Pn}{\mathcal{P}}
\newcommand{\Span}{\operatorname{L}}
\newcommand{\rg}{\operatorname{rg}}
\newcommand{\dimn}{\operatorname{dim}}

\begin{document}

\begin{center}
	{\LARGE\bfseries Álgebra Lineal}\\[0.4em]
	{\Large Tema 1: Espacios Vectoriales}
\end{center}

\tableofcontents
\newpage

\section{Espacio vectorial}

\subsection{Definición}

Un espacio vectorial sobre un cuerpo $\K$ es un conjunto no vacío $E$ cuyos elementos se denominan vectores, dotado de una suma y de un producto por escalares.

El conjunto $(E,+)$ debe ser un grupo conmutativo. Para todo $x,y,z\in E$:

\[
x+y\in E
\]

\[
(x+y)+z=x+(y+z)
\]

\[
x+0=x
\]

\[
\forall x\in E,\qquad \exists -x\in E
\quad\text{tal que}\quad
x+(-x)=0
\]

\[
x+y=y+x
\]

Además, para todo $x,y\in E$ y $\gamma,\mu\in\K$:

\[
(\gamma+\mu)x=\gamma x+\mu x
\]

\[
\gamma(x+y)=\gamma x+\gamma y
\]

\[
\gamma(\mu x)=(\gamma\mu)x
\]

\[
1x=x
\]

\subsection{Propiedades}

En todo espacio vectorial se verifica:

\[
0x=0
\]

\[
\gamma 0=0
\]

\[
\gamma x=0
\quad\Longrightarrow\quad
\gamma=0
\quad\text{o}\quad
x=0
\]

\[
(-\gamma)x=-(\gamma x)=\gamma(-x)
\]

Si $u\neq0$:

\[
\gamma u=\mu u
\quad\Longrightarrow\quad
\gamma=\mu
\]

Si $\gamma\neq0$:

\[
\gamma u=\gamma v
\quad\Longrightarrow\quad
u=v
\]

\subsection{Ejemplos}

Son espacios vectoriales:

\begin{itemize}
	\item $\R^n$.
	\item El conjunto $\Pn_n(x)$ de polinomios de grado menor o igual que $n$.
	\item El conjunto $M_{m\times n}(\R)$ de matrices $m\times n$ con coeficientes reales.
	\item El conjunto de soluciones de un sistema homogéneo de ecuaciones lineales.
\end{itemize}

\section{Subespacios vectoriales}

Un subconjunto no vacío $S\subseteq E$ es un subespacio vectorial de $E$ si, con las operaciones heredadas de $E$, tiene estructura de espacio vectorial.

Para comprobarlo basta verificar:

\[
S\neq\emptyset
\]

\[
\forall x,y\in S,
\qquad
x+y\in S
\]

\[
\forall\alpha\in\R,\ \forall x\in S,
\qquad
\alpha x\in S
\]

De forma equivalente:

\[
\boxed{
S\leq E
\iff
S\neq\emptyset
\ \text{y}\
\alpha x+\beta y\in S
\quad
\forall x,y\in S,\ \forall\alpha,\beta\in\R
}
\]

\subsection{Ejemplos}

El conjunto

\[
S=
\left\{
(x,y,z,t)\in\R^4:
x+2y-z+3t=0
\right\}
\]

es un subespacio vectorial de $\R^4$.

También lo es:

\[
T=
\left\{
(x,y,z)\in\R^3:
\begin{array}{rcl}
x-3y+z&=&0\\
2x-z&=&0
\end{array}
\right\}
\]

En cambio:

\[
M=
\left\{
(x,y,z)\in\R^3:
xy=0
\right\}
\]

no es un subespacio vectorial.

Tampoco lo es:

\[
S=
\left\{
(x,y,z,t)\in\R^4:
x+2y-z+3t=1
\right\}
\]

porque no contiene al vector nulo.

\section{Combinación lineal y subespacio engendrado}

\subsection{Combinación lineal}

Un vector $x\in E$ es combinación lineal de los vectores

\[
S=\{v_1,v_2,\ldots,v_n\}
\]

si existen escalares $\alpha_1,\ldots,\alpha_n\in\K$ tales que:

\[
x=
\alpha_1v_1+\alpha_2v_2+\cdots+\alpha_nv_n
\]

Por ejemplo:

\[
(1,-1,3)
=
(1,0,0)-(0,1,0)+3(0,0,1)
\]

\subsection{Subespacio engendrado}

El subespacio engendrado por $S=\{v_1,\ldots,v_n\}$ es:

\[
\Span(S)
=
\left\{
\alpha_1v_1+\cdots+\alpha_nv_n:
\alpha_1,\ldots,\alpha_n\in\K
\right\}
\]

$\Span(S)$ es el menor subespacio vectorial que contiene a $S$.

Por ejemplo:

\[
S=
\left\{
(\alpha,\beta,2\alpha+\beta,-\alpha+3\beta):
\alpha,\beta\in\R
\right\}
\]

puede escribirse como:

\[
S=
\Span\left\{
(1,0,2,-1),
(0,1,1,3)
\right\}
\]

y su expresión implícita es:

\[
S=
\left\{
(x,y,z,t)\in\R^4:
z=2x+y,\quad
t=-x+3y
\right\}
\]

\subsection{Sistema generador}

$S=\{v_1,\ldots,v_n\}$ es un sistema generador de un espacio vectorial $V$ si todo vector $v\in V$ puede expresarse como combinación lineal de los vectores de $S$:

\[
\forall v\in V,
\qquad
v=
\alpha_1v_1+\cdots+\alpha_nv_n
\]

para ciertos $\alpha_1,\ldots,\alpha_n\in\K$.

\section{Independencia lineal}

Los vectores $v_1,\ldots,v_n\in E$ son linealmente independientes si:

\[
\alpha_1v_1+\cdots+\alpha_nv_n=0
\]

implica necesariamente:

\[
\alpha_1=\alpha_2=\cdots=\alpha_n=0
\]

Un conjunto de vectores linealmente independientes se denomina sistema libre.

Si existe una combinación lineal nula con algún coeficiente no nulo, el sistema es linealmente dependiente o ligado.

\subsection{Propiedades}

\begin{itemize}
	\item Si $v\neq0$, entonces $\{v\}$ es un sistema libre.
	\item Todo subconjunto de un sistema libre es libre.
	\item Todo conjunto que contiene a un sistema ligado es ligado.
	\item Todo sistema que contiene al vector $0$ es ligado.
	\item Un sistema es ligado si alguno de sus vectores es combinación lineal de los restantes.
	\item Un sistema es libre si ninguno de sus vectores puede expresarse como combinación lineal de los restantes.
\end{itemize}

Para un conjunto de $n$ vectores de $\R^n$, la independencia lineal puede comprobarse mediante:

\[
\rg(A)=n
\]

o, si $A$ es cuadrada:

\[
\det(A)\neq0
\]

\subsection{Teorema fundamental de la independencia lineal}

Si

\[
S=\{v_1,\ldots,v_n\}
\]

es un sistema generador de $V$ y

\[
T=\{w_1,\ldots,w_p\}
\]

es un sistema linealmente independiente de vectores de $V$, entonces:

\[
\boxed{p\leq n}
\]

\section{Base y dimensión}

\subsection{Base}

Sea $V$ un espacio vectorial de dimensión finita. Un conjunto

\[
B=\{e_1,\ldots,e_n\}
\]

es una base de $V$ si sus vectores son linealmente independientes y forman un sistema generador.

Equivalentemente, todo vector $v\in V$ se expresa de forma única como combinación lineal de los vectores de $B$.

Bases canónicas habituales:

\[
B_c(\R^n)
=
\{e_1,\ldots,e_n\}
\]

\[
B_c(\Pn_n(x))
=
\{1,x,x^2,\ldots,x^n\}
\]

y en $M_{2\times2}(\R)$:

\[
\left\{
\begin{pmatrix}1&0\\0&0\end{pmatrix},
\begin{pmatrix}0&1\\0&0\end{pmatrix},
\begin{pmatrix}0&0\\1&0\end{pmatrix},
\begin{pmatrix}0&0\\0&1\end{pmatrix}
\right\}
\]

\subsection{Teorema de ampliación de la base}

Si $S$ es un subespacio vectorial de $E$ y $B_S$ es una base de $S$, existen vectores

\[
v_1,\ldots,v_r\in E
\]

tales que:

\[
B_S\cup\{v_1,\ldots,v_r\}
\]

es una base de $E$.

El subespacio generado por los vectores añadidos se denomina subespacio suplementario de $S$.

\subsection{Dimensión}

La dimensión de un espacio vectorial de dimensión finita es el número de elementos de cualquiera de sus bases.

\[
\boxed{
\dimn V
=
\text{número de vectores de una base de }V
}
\]

Además:

\begin{itemize}
	\item $\dimn V$ es el máximo número de vectores linealmente independientes de $V$.
	\item $\dimn V$ es el mínimo número de vectores necesarios para generar $V$.
	\item Si $\dimn V=n$, todo conjunto de $n$ vectores linealmente independientes es una base.
	\item Si $S\subseteq T$, entonces $\dimn S\leq\dimn T$.
	\item Si $S\subseteq T$ y $\dimn S=\dimn T$, entonces $S=T$.
\end{itemize}

En $\R^n$, la dimensión de un subespacio coincide con el número de parámetros libres de su expresión paramétrica.

\subsection{Rango de un sistema de vectores}

El rango de un sistema de vectores $S$ es la dimensión del subespacio engendrado:

\[
\boxed{
\rg(S)=\dimn\Span(S)
}
\]

Es decir, es el máximo número de vectores linealmente independientes del sistema.

\section{Coordenadas y cambio de base}

\subsection{Coordenadas de un vector}

Sea:

\[
B=\{e_1,\ldots,e_n\}
\]

una base de $V$. Si:

\[
x=
\alpha_1e_1+\cdots+\alpha_ne_n
\]

entonces:

\[
[x]_B=
\begin{pmatrix}
\alpha_1\\
\vdots\\
\alpha_n
\end{pmatrix}
\]

son las coordenadas de $x$ en la base $B$.

Las coordenadas de un vector en una base son únicas.

\subsection{Matriz de cambio de base}

Sean $B$ y $B'$ dos bases de un espacio vectorial $V$.

La matriz de cambio de base de $B$ a $B'$ se construye colocando como columnas las coordenadas de los vectores de $B$ expresados en la base $B'$:

\[
\boxed{
M_{BB'}
=
\left(
[b_1]_{B'}
\ \cdots\
[b_n]_{B'}
\right)
}
\]

y transforma coordenadas según:

\[
\boxed{
[x]_{B'}
=
M_{BB'}[x]_B
}
\]

La matriz inversa realiza el cambio contrario:

\[
M_{B'B}
=
M_{BB'}^{-1}
\]

Toda matriz de cambio de base es cuadrada y regular.

Si $A$, $B$ y $C$ son tres bases:

\[
\boxed{
M_{AB}
=
M_{CB}M_{AC}
}
\]

\section{Suma e intersección de subespacios}

\subsection{Intersección}

Dados dos subespacios $S_1,S_2\leq E$:

\[
S_1\cap S_2
=
\{x\in E:
x\in S_1
\text{ y }
x\in S_2
\}
\]

La intersección de subespacios es siempre un subespacio vectorial.

\subsection{Suma}

La suma de dos subespacios se define como:

\[
S_1+S_2
=
\left\{
u_1+u_2:
u_1\in S_1,\;
u_2\in S_2
\right\}
\]

También es un subespacio vectorial.

\subsection{Fórmula de Grassmann}

\[
\boxed{
\dimn(S_1\cap S_2)
+
\dimn(S_1+S_2)
=
\dimn S_1
+
\dimn S_2
}
\]

\subsection{Suma directa}

La suma de dos subespacios $S$ y $T$ es directa si:

\[
S\cap T=\{0\}
\]

y se escribe:

\[
S\oplus T
\]

En ese caso:

\[
\dimn(S\oplus T)
=
\dimn S+\dimn T
\]

y todo vector de $S+T$ se expresa de forma única como suma de un vector de $S$ y otro de $T$.

Si además:

\[
S\oplus T=E
\]

entonces $S$ y $T$ son subespacios suplementarios.

\section{Espacio vectorial producto}

Sean $E$ y $F$ espacios vectoriales sobre el mismo cuerpo $\K$.

El producto cartesiano:

\[
E\times F
\]

con las operaciones:

\[
(u_1,u_2)+(v_1,v_2)
=
(u_1+v_1,u_2+v_2)
\]

\[
\gamma(u,v)
=
(\gamma u,\gamma v)
\]

es un espacio vectorial.

Si:

\[
\dimn E=n,
\qquad
\dimn F=m
\]

entonces:

\[
\boxed{
\dimn(E\times F)=n+m
}
\]

Si $\{e_1,\ldots,e_n\}$ es una base de $E$ y $\{f_1,\ldots,f_m\}$ una base de $F$, entonces una base de $E\times F$ es:

\[
\left\{
(e_1,0),\ldots,(e_n,0),
(0,f_1),\ldots,(0,f_m)
\right\}
\]

\section{Espacio vectorial cociente}

Sea $V$ un subespacio de $E$. Se define la relación:

\[
x\sim y
\iff
x-y\in V
\]

Las clases de equivalencia forman el espacio cociente:

\[
E/V
=
\{[x]:
x\in E\}
\]

con:

\[
[x]+[y]=[x+y]
\]

\[
\gamma[x]=[\gamma x]
\]

Si:

\[
\dimn E=n,
\qquad
\dimn V=m
\]

entonces:

\[
\boxed{
\dimn(E/V)=n-m
}
\]

\end{document}
